% =============================================================
% Capítulo 5. Implementación
% =============================================================
% Realización en código de la propuesta: tecnologías, estructura del
% proyecto, integración con Forge, paralelización y control de versiones.
% =============================================================

\chapter{Implementación}
\label{cap:implementacion}

Este capítulo describe cómo se ha realizado en código la propuesta del capítulo anterior, esto es, las tecnologías empleadas, la estructura del proyecto y los aspectos técnicos que hacen viable la evaluación por simulación.

\section{Tecnologías y estructura del proyecto}

El sistema se ha implementado principalmente en Python, que concentra el algoritmo genético y la lógica de orquestación, e integra el simulador Forge, escrito en Java, como motor externo de evaluación. El catálogo de cartas se obtiene de la API pública de Scryfall.

El código se organiza en cuatro módulos, cada uno responsable de una etapa del flujo descrito en la sección~\ref{sec:pipeline}:

\begin{itemize}
    \item \texttt{obtener\_cartas\_mtg.py} descarga las cartas de Scryfall, restringidas a las expansiones del formato Standard, y genera dos archivos: un catálogo (\texttt{card\_catalog.json}), que asigna a cada carta un identificador entero y almacena sus atributos, y unos índices (\texttt{card\_indices.json}), que agrupan las cartas por tipo y color para acelerar las búsquedas.
    \item \texttt{generador\_mazos\_mtg.py} produce la población inicial de mazos con las restricciones del formato.
    \item \texttt{algoritmo\_genetico\_mtg.py} es el núcleo del sistema: su clase \texttt{MTG\allowbreak Genetic\allowbreak Algorithm} encapsula el ciclo evolutivo completo: inicialización, evaluación, selección, cruce, mutación, elitismo y persistencia.
    \item \texttt{mtg\_main.py} actúa como orquestador: conecta las tres etapas y analiza el hardware disponible antes de cada ejecución.
\end{itemize}

El proyecto se mantiene bajo control de versiones con Git, en un repositorio público\footnote{Repositorio del proyecto: \url{https://github.com/JGasBul/TFGJoseGascoBule}}. El repositorio incluye un fichero README con las instrucciones para instalar las dependencias, obtener el catálogo de cartas y lanzar las pruebas, de modo que cualquier interesado pueda reproducir los experimentos.

\section{Evaluación paralela e integración con Forge}

La medida del fitness exige lanzar Forge como proceso externo. Cada enfrentamiento invoca el modo de simulación del motor (\texttt{forge sim}), que disputa tres partidas entre dos mazos con jugador inicial aleatorio. El sistema interpreta la salida del simulador, cuenta las victorias y determina el ganador.

El coste de simular partidas reales obliga a paralelizar la evaluación. Los enfrentamientos de una misma generación se reparten entre varios procesos mediante un \texttt{ProcessPoolExecutor}, cuyo número de trabajadores fija la clase \texttt{HardwareAnalyzer} en función de los núcleos de CPU, la memoria y la velocidad de disco de la máquina. De este modo, el mismo código se ejecuta en un ordenador portátil o en un servidor de muchos núcleos sin reconfiguración manual. Un límite de tiempo adaptativo evita que una partida bloqueada detenga la ejecución: si se acumulan interrupciones por tiempo, el límite se eleva de forma automática. En las ejecuciones reales, la tasa de partidas fallidas se mantuvo por debajo del 0.1\,\%, lo que confirma la estabilidad de la integración a lo largo de decenas de horas de cómputo.

\section{Persistencia y reproducibilidad}

Al término de cada generación, el sistema guarda un punto de control con el estado de la población, el Hall of Fame y las estadísticas de la ejecución. Esta persistencia cumple dos funciones: permite reanudar una ejecución interrumpida sin repetir el cómputo ya realizado y, sobre todo, hace posible reconstruir a posteriori la trayectoria de cada experimento, lo que sustenta el análisis del capítulo siguiente. Para disputar cada partida, el fenotipo de un mazo se traduce al archivo de mazo (\texttt{.dck}) que Forge lee, de modo que cada individuo evaluado queda también almacenado en un formato reutilizable.
